The Moon and planets are shown in the projected sky, if they are above the
horizon. The Milky Way and other diffuse objects are not shown. In the
planetarium's projection, azimuth starts at the North Cardinal Point, which
is assigned 0 degrees; the East Cardinal Point has an azimuth of 90 degrees,
and so on.

\begin{description}
\item[Question 1 (4 points)]~
  \begin{description}
  \item[(a)] In which constellation is the Moon located? Give the Latin name
    or the IAU abbreviation.
  \item[(b)] Given that, for this projected sky, the Sun has set 1 hour and
    26 minutes ago, identify the current month.
  \end{description}
\item[Question 2 (4 points)] Which of the following stars are circumpolar?
  Circle your answers.

  \begin{center}
  \begin{tabular}{cccc}
  $\alpha$ Car & $\gamma$ Cru & $\beta$ Hyi & $\beta$ Pav \\
  $\beta$ Crv & $\beta$ Car & $\alpha$ Eri & $\upsilon$ Oct \\
  $\alpha$ Sco & $\beta$ Lup & $\gamma$ Vel & $\alpha$ Cnc \\
  \end{tabular}
  \end{center}
\item[Question 3 (2 points)] A certain artificial satellite was seen passing
  through this sky with the following coordinates:
  \begin{itemize}
  \item At rise: Azimuth: $98^\circ 47'$; Height: $7^\circ 31'$
  \item At maximum brightness: Azimuth: $2^\circ 27'$; Height:
    $54^\circ 47'$
  \item At disappearance: Azimuth: $311^\circ 19'$; Height: $26^\circ 47'$
  \end{itemize}
  Using the azimuth references given above, write the Bayer designation of
  one bright star, for each case, very close to the position of the
  satellite, when it is at:
  \begin{description}
  \item[(a)] Maximum brightness
  \item[(b)] Disappearance
  \end{description}
\end{description}
